% answers.tex -- answers and worked solutions for booklet.tex (Module 6 MC booklet), same
% layout (quiz.sty). One practiceans per quizq, in the same order (72).
% Every number: checks/m6-mc-booklet/numbers.py
% Build: TEXINPUTS=.: xelatex answers.tex (twice)
\documentclass[12pt]{article}
\usepackage{quiz}
\hypersetup{pdftitle={Multiple Choice Booklet - Module 6 - Electromagnetism - Answers}}
% links between the two PDF files are not possible, so question headings are plain text here
\renewcommand\hyperlink[2]{#2}
\newcommand\anstopic[1]{\par\Needspace{8\baselineskip}\vspace{4mm}%
  \begin{center}{\headfont\bfseries\fontsize{14}{18}\selectfont #1\par}\end{center}\vspace{2pt}}

\begin{document}
\quiztitle{Year 12 Physics}{Multiple Choice Booklet - Module 6 - Electromagnetism: Answers}

\begin{center}
  \renewcommand\arraystretch{1.45}
  \begin{tabular}{|l|*{12}{>{\centering\arraybackslash}p{7mm}|}}
    \hline
    \textbf{Q} & 1 & 2 & 3 & 4 & 5 & 6 & 7 & 8 & 9 & 10 & 11 & 12\\ \hline
    \textbf{Ans} & C & B & D & B & D & C & D & A & B & B & C & A\\ \hline\hline
    \textbf{Q} & 13 & 14 & 15 & 16 & 17 & 18 & 19 & 20 & 21 & 22 & 23 & 24\\ \hline
    \textbf{Ans} & C & C & A & A & B & D & C & C & B & A & B & D\\ \hline\hline
    \textbf{Q} & 25 & 26 & 27 & 28 & 29 & 30 & 31 & 32 & 33 & 34 & 35 & 36\\ \hline
    \textbf{Ans} & C & A & A & B & D & C & A & C & B & D & A & D\\ \hline\hline
    \textbf{Q} & 37 & 38 & 39 & 40 & 41 & 42 & 43 & 44 & 45 & 46 & 47 & 48\\ \hline
    \textbf{Ans} & B & D & C & A & D & D & C & D & A & B & C & B\\ \hline\hline
    \textbf{Q} & 49 & 50 & 51 & 52 & 53 & 54 & 55 & 56 & 57 & 58 & 59 & 60\\ \hline
    \textbf{Ans} & C & B & D & B & A & A & D & B & C & C & A & B\\ \hline\hline
    \textbf{Q} & 61 & 62 & 63 & 64 & 65 & 66 & 67 & 68 & 69 & 70 & 71 & 72\\ \hline
    \textbf{Ans} & A & C & D & A & D & C & B & A & B & C & D & B\\ \hline
  \end{tabular}
\end{center}
Topic 1: Questions 1--18. \quad Topic 2: Questions 19--36. \quad Topic 3: Questions 37--54. \quad
Topic 4: Questions 55--72.

Directions: right-hand palm rule for a \emph{positive} charge (or conventional current): fingers along
$B$, thumb along $v$ (or $I$), the palm pushes in the direction of the force. For an electron, the
force is in the opposite direction.

\begin{practiceanswers}

% =====================================================================
\anstopic{Topic 1 - Charged Particles in Electric and Magnetic Fields}

% ---- Q1
\begin{practiceans}
  \correct{C}
  $F = qvB\sin\theta$. An electron has the same size of charge as a proton ($e$). At the same speed in
  the same field, the force has the same size (it acts in the opposite direction).
  \begin{itemize}
    \item \xmark\ A and B: $F \propto B$ and $F \propto v$, so doubling either doubles the force.
    \item \xmark\ D: an alpha particle has charge $2e$, so the force doubles.
  \end{itemize}
\end{practiceans}

% ---- Q2
\Needspace{9\baselineskip}
\begin{practiceans}
  \correct{B}
  An electric force acts along the field and can do work on the charge ($W = qEd$), so it can change
  the kinetic energy. A magnetic force is always at \ang{90} to the velocity, so it does no work: it
  changes only the direction of motion, not the speed.
  \begin{itemize}
    \item \xmark\ A and C: a magnetic field cannot change the speed or kinetic energy.
    \item \xmark\ D: an electric field can.
  \end{itemize}
\end{practiceans}

% ---- Q3
\Needspace{9\baselineskip}
\begin{practiceans}
  \correct{D}
  $E = \dfrac{V}{d} = \dfrac{250}{\num{2.0e-3}} = \qty{1.25e5}{V.m^{-1}}$.
  \begin{itemize}
    \item \xmark\ A: this is $\frac{d}{V}$. \quad \xmark\ B: this is $V \times d$.
    \item \xmark\ C: uses $d$ in millimetres instead of metres.
  \end{itemize}
\end{practiceans}

% ---- Q4
\Needspace{9\baselineskip}
\begin{practiceans}
  \correct{B}
  The magnetic force provides the centripetal force: $qvB = \dfrac{mv^2}{r}$, so
  \[ r = \frac{mv}{qB} = \frac{(\num{1.673e-27})(\num{3.0e6})}{(\num{1.602e-19})(0.20)} = \qty{0.16}{m} \]
  \begin{itemize}
    \item \xmark\ A: uses the mass of an electron. \quad \xmark\ C: this is $\frac{qB}{mv}$ (upside down).
    \item \xmark\ D: uses $v^2$ instead of $v$.
  \end{itemize}
\end{practiceans}

% ---- Q5
\Needspace{9\baselineskip}
\begin{practiceans}
  \correct{D}
  The path curves upwards, so the force on the electron is up the page (towards the centre of the
  curve). For a \emph{positive} charge moving right, an upward force needs $B$ into the page. The
  electron is negative, so its force is reversed: $B$ must be out of the page.
  \begin{itemize}
    \item \xmark\ A and B: a field up or down the page would give a force into or out of the page, so
          the electron would leave the plane of the page.
    \item \xmark\ C: a field into the page would push the electron down.
  \end{itemize}
\end{practiceans}

% ---- Q6
\Needspace{9\baselineskip}
\begin{practiceans}
  \correct{C}
  At speed $v$ the forces balance: $qE = qvB$. The magnetic force is proportional to the speed, but the
  electric force is not. At $2v$ the magnetic force doubles while the electric force stays the same, so
  the net force is in the direction of the magnetic force, and the ion is deflected that way.
  \begin{itemize}
    \item \xmark\ A: the forces balance only at $v = \frac{E}{B}$.
    \item \xmark\ B: the magnetic force is now the larger force.
    \item \xmark\ D: the electric force $qE$ does not depend on speed.
  \end{itemize}
\end{practiceans}

% ---- Q7
\Needspace{9\baselineskip}
\begin{practiceans}
  \correct{D}
  $r = \dfrac{mv}{qB}$. Doubling $v$ doubles $r$; halving $B$ doubles $r$ again. So $r$ becomes $4r$.
  \begin{itemize}
    \item \xmark\ A: the changes are inverted. \quad \xmark\ B: the two changes do not cancel; both
          increase $r$.
    \item \xmark\ C: includes only one of the two changes.
  \end{itemize}
\end{practiceans}

% ---- Q8
\Needspace{9\baselineskip}
\begin{practiceans}
  \correct{A}
  $F = qvB\sin\theta$, with the same $q$ and $B$ each time, so compare $v\sin\theta$:
  A: $50 \times \sin 90^\circ = 50$. B: $80 \times \sin 30^\circ = 40$. C: $100 \times \sin 0^\circ = 0$.
  D: $60 \times \sin 45^\circ = 42$. The largest is A.
  \begin{itemize}
    \item \xmark\ B and D: only the component of $v$ at right angles to $B$ counts.
    \item \xmark\ C: a charge moving parallel to the field feels no magnetic force.
  \end{itemize}
\end{practiceans}

% ---- Q9
\Needspace{9\baselineskip}
\begin{practiceans}
  \correct{B}
  A proton has the opposite charge, so it is deflected to the other side of $P$. It has the same size
  of charge, so the same force, but about 1840 times the mass. Its acceleration, and so its deflection,
  is much smaller. It hits just to the right of $P$, at $B$.
  \begin{itemize}
    \item \xmark\ A: a positive charge is pushed the opposite way to an electron.
    \item \xmark\ C: an equal deflection would need an equal mass.
    \item \xmark\ D: the proton's deflection is smaller, not larger.
  \end{itemize}
\end{practiceans}

% ---- Q10
\Needspace{9\baselineskip}
\begin{practiceans}
  \correct{B}
  $E = \dfrac{V}{d}$. Doubling both $V$ and $d$ leaves $E$ unchanged, so the force $qE$ is the same.
  Over the same distance $XY$, the work done (gain in kinetic energy) is the same: $W$.
  \begin{itemize}
    \item \xmark\ A: only the distance between the plates was considered.
    \item \xmark\ C and D: only the voltage was considered.
  \end{itemize}
\end{practiceans}

% ---- Q11
\Needspace{9\baselineskip}
\begin{practiceans}
  \correct{C}
  The proton is deflected down, so the electric force on it is down. A positive charge is pushed along
  the field, so $E$ points towards the bottom of the page. The magnetic force must be up. For a positive
  charge moving right, $B$ into the page gives an upward force (palm rule).
  \begin{itemize}
    \item \xmark\ A and B: an upward field would push the proton up, but it was deflected down.
    \item \xmark\ D: $B$ out of the page gives a downward force, adding to the electric force.
  \end{itemize}
\end{practiceans}

% ---- Q12
\Needspace{9\baselineskip}
\begin{practiceans}
  \correct{A}
  $E = \dfrac{V}{d} = \dfrac{800}{\num{2.0e-3}} = \qty{4.0e5}{V.m^{-1}}$. For no deflection,
  $qE = qvB$, so $B = \dfrac{E}{v} = \dfrac{\num{4.0e5}}{\num{4.0e6}} = \qty{0.10}{T}$.
  The top plate is positive, so the electric force on the electrons is up; the magnetic force must be
  down. For an electron moving right, $B$ into the page gives a downward force.
  \begin{itemize}
    \item \xmark\ B: this field would push the electrons up, adding to the electric force.
    \item \xmark\ C and D: these use $\frac{v}{E}$ instead of $\frac{E}{v}$.
  \end{itemize}
\end{practiceans}

% ---- Q13
\Needspace{9\baselineskip}
\begin{practiceans}
  \correct{C}
  The top plate is positive, so the electric force on the electron is up. The magnetic force must be
  down. With $B$ out of the page, a positive charge moving left would feel an upward force, so an
  electron moving left feels a downward force. The forces balance, so the velocity is constant.
  \begin{itemize}
    \item \xmark\ A and B: a velocity up or down gives a sideways magnetic force, which cannot balance
          the vertical electric force.
    \item \xmark\ D: an electron moving right would feel an upward magnetic force.
  \end{itemize}
\end{practiceans}

% ---- Q14
\Needspace{9\baselineskip}
\begin{practiceans}
  \correct{C}
  Work done $=$ gain in kinetic energy: $qV = \frac{1}{2}mv^2$, so
  \[ v = \sqrt{\frac{2qV}{m}} = \sqrt{\frac{2(\num{1.602e-19})(2000)}{\num{9.109e-31}}} = \qty{2.7e7}{m.s^{-1}} \]
  \begin{itemize}
    \item \xmark\ A: uses the proton's mass. \quad \xmark\ B: forgets the 2.
    \item \xmark\ D: forgets the square root (this is $v^2$).
  \end{itemize}
\end{practiceans}

% ---- Q15
\Needspace{9\baselineskip}
\begin{practiceans}
  \correct{A}
  In the gravitational field, the acceleration is $g$, whatever the charge: the path and landing point
  $X$ do not change. In the electric field, $a = \frac{qE}{m}$ doubles, so the particle falls faster and
  lands sooner, closer to the start than $Y$.
  \begin{itemize}
    \item \xmark\ B and D: the motion in a gravitational field does not depend on charge.
    \item \xmark\ C: the electric acceleration increases, so the particle lands closer, not further.
  \end{itemize}
\end{practiceans}

% ---- Q16
\Needspace{9\baselineskip}
\begin{practiceans}
  \correct{A}
  The horizontal speed and distance are the same, so the time is the same. The drop $\frac{1}{2}at^2$
  halves, so the downward acceleration is $\frac{g}{2}$. The net force is $\frac{mg}{2}$ down:
  $mg - qE = \dfrac{mg}{2}$, so $E = \dfrac{mg}{2q}$ (directed upwards).
  \begin{itemize}
    \item \xmark\ B: $E = \frac{mg}{q}$ balances gravity, so the particle would not fall at all.
    \item \xmark\ C and D: these give a net upward force, so the particle would rise.
  \end{itemize}
\end{practiceans}

% ---- Q17
\Needspace{9\baselineskip}
\begin{practiceans}
  \correct{B}
  Electrons moving in the same direction as $I$ form a conventional current in the \emph{opposite}
  direction. Antiparallel currents repel, so the beam is pushed away from the wire. (Check: above the
  wire, its magnetic field points out of the page; the palm rule for a positive charge moving with $I$
  gives a force towards the wire, so for an electron the force is away from the wire.)
  \begin{itemize}
    \item \xmark\ A: this would happen if the conventional currents were in the same direction.
    \item \xmark\ C: the force is in the plane of the page. \quad \xmark\ D: the wire's magnetic
          field does act on moving charges.
  \end{itemize}
\end{practiceans}

% ---- Q18
\Needspace{9\baselineskip}
\begin{practiceans}
  \correct{D}
  For the first half there is no field, so the mass falls straight down. Then a constant horizontal
  force acts, towards the negative (right) plate, as well as gravity. At that moment the velocity is
  vertical and the net force is at an angle to it, so the path curves to the right, like a projectile.
  \begin{itemize}
    \item \xmark\ A: a straight diagonal path from the start would need the field to be on from the start.
    \item \xmark\ B: the velocity cannot change direction suddenly (a sharp corner).
    \item \xmark\ C: the path should be vertical for the first half.
  \end{itemize}
\end{practiceans}

% =====================================================================
\anstopic{Topic 2 - The Motor Effect}

% ---- Q19
\begin{practiceans}
  \correct{C}
  Between the N pole of the left magnet and the S pole of the right magnet, $B$ points to the right.
  Palm rule: fingers right ($B$), thumb out of the page ($I$): the palm pushes up the page.
  \begin{itemize}
    \item \xmark\ A and B: the force is at right angles to $B$, so it cannot be left or right.
    \item \xmark\ D: this is the reverse direction.
  \end{itemize}
\end{practiceans}

% ---- Q20
\Needspace{9\baselineskip}
\begin{practiceans}
  \correct{C}
  The top and bottom sides are parallel to $B$, so they feel no force. On the left side (current up,
  $B$ right) the force is into the page; on the right side (current down) it is out of the page. These
  two forces turn the loop about the vertical axis $WX$, with the left side moving into the page.
  \begin{itemize}
    \item \xmark\ A and B: turning about $YZ$ would need forces on the top and bottom sides.
    \item \xmark\ D: this is the wrong direction.
  \end{itemize}
\end{practiceans}

% ---- Q21
\Needspace{9\baselineskip}
\begin{practiceans}
  \correct{B}
  $\dfrac{F}{l} = \dfrac{\mu_0 I_1 I_2}{2\pi r} = \dfrac{(4\pi\times10^{-7})(2.0)(3.0)}{2\pi(0.50)}
  = \qty{2.4e-6}{N.m^{-1}}$.
  \begin{itemize}
    \item \xmark\ A: multiplies by $r$. \quad \xmark\ C: uses $r^2$.
    \item \xmark\ D: leaves out the $2\pi$.
  \end{itemize}
\end{practiceans}

% ---- Q22
\Needspace{9\baselineskip}
\begin{practiceans}
  \correct{A}
  $\dfrac{F}{l} = \dfrac{\mu_0 I_1 I_2}{2\pi r}$, so $r = \dfrac{\mu_0 I_1 I_2 l}{2\pi F}$.
  \begin{itemize}
    \item \xmark\ B: this is upside down. \quad \xmark\ C: $l$ is on the wrong side.
    \item \xmark\ D: there is no square root.
  \end{itemize}
\end{practiceans}

% ---- Q23
\Needspace{9\baselineskip}
\begin{practiceans}
  \correct{B}
  $F = lIB\sin\theta = 0.080 \times 3.0 \times 0.50 \times \sin 40^\circ = \qty{0.077}{N}$.
  \begin{itemize}
    \item \xmark\ A: the conductor is not parallel to the field. \quad \xmark\ C: uses $\cos 40^\circ$.
    \item \xmark\ D: ignores the angle.
  \end{itemize}
\end{practiceans}

% ---- Q24
\Needspace{9\baselineskip}
\begin{practiceans}
  \correct{D}
  The currents are in opposite directions, so the wires repel: $X$ pushes $Y$ down onto the balance and
  the reading is more than the mass of $Y$. A larger current in $X$ gives a larger repulsion, so the
  reading increases.
  \begin{itemize}
    \item \xmark\ A: reversing $Y$ makes the currents parallel, so they attract and the reading decreases.
    \item \xmark\ B: there is a magnetic force on $Y$.
    \item \xmark\ C: reversing both keeps them opposite, so the reading does not change.
  \end{itemize}
\end{practiceans}

% ---- Q25
\Needspace{9\baselineskip}
\begin{practiceans}
  \correct{C}
  $F = lIB\sin\theta$. At first $\theta = 90^\circ$. The force halves when $\sin\theta = 0.5$, so
  $\theta = 30^\circ$ to the field. The rod has turned through $90^\circ - 30^\circ = 60^\circ$.
  \begin{itemize}
    \item \xmark\ A: turning \ang{30} leaves $\theta = 60^\circ$, and $F = 0.87F$.
    \item \xmark\ B: gives $0.71F$. \quad \xmark\ D: the rod would be parallel to the field, so $F = 0$.
  \end{itemize}
\end{practiceans}

% ---- Q26
\Needspace{9\baselineskip}
\begin{practiceans}
  \correct{A}
  Palm rule: fingers up ($B$), thumb east ($I$): the palm faces south, so the force is to the south.
  \begin{itemize}
    \item \xmark\ B: this is the reverse direction.
    \item \xmark\ C and D: the force is at right angles to $B$, which is vertical, so it cannot be vertical.
  \end{itemize}
\end{practiceans}

% ---- Q27
\Needspace{9\baselineskip}
\begin{practiceans}
  \correct{A}
  $Q$ and $R$ carry currents in the same direction, so they attract: $Q$ pulls $R$ to the left with force
  $F$. $P$ and $R$ carry opposite currents, so they repel: $P$ pushes $R$ to the right with $\frac{F}{2}$
  (twice the distance). Net force $= F - \frac{F}{2} = \frac{F}{2}$ to the left.
  \begin{itemize}
    \item \xmark\ B: the wrong direction. \quad \xmark\ C and D: these add the two forces, but they act
          in opposite directions.
  \end{itemize}
\end{practiceans}

% ---- Q28
\Needspace{9\baselineskip}
\begin{practiceans}
  \correct{B}
  $F \propto \dfrac{I_1 I_2}{r}$: the currents give $3 \times 3 = 9$, the distance divides by 3, so the
  force is multiplied by 3.
  \begin{itemize}
    \item \xmark\ A: the changes do not cancel. \quad \xmark\ C: forgets the larger distance.
    \item \xmark\ D: forgets the larger currents.
  \end{itemize}
\end{practiceans}

% ---- Q29
\Needspace{9\baselineskip}
\begin{practiceans}
  \correct{D}
  $I = \dfrac{F}{lB\sin\theta} = \dfrac{0.024}{0.30 \times 0.040 \times \sin 30^\circ} = \qty{4.0}{A}$.
  \begin{itemize}
    \item \xmark\ A: puts $\sin\theta$ on the top. \quad \xmark\ B: ignores the angle.
    \item \xmark\ C: uses $\cos 30^\circ$.
  \end{itemize}
\end{practiceans}

% ---- Q30
\Needspace{9\baselineskip}
\begin{practiceans}
  \correct{C}
  $F = \dfrac{\mu_0 I_1 I_2 l}{2\pi r} = \dfrac{(2\times10^{-7})(30)(40)(1.50)}{\num{3.0e-3}} = \qty{0.12}{N}$.
  Each trolley feels this force (Newton's third law), so $m = \dfrac{F}{a} = \dfrac{0.12}{0.050} = \qty{2.4}{kg}$.
  \begin{itemize}
    \item \xmark\ A: shares the force between the trolleys; each one feels the whole force.
    \item \xmark\ B: forgets the length $l$. \quad \xmark\ D: doubles the force.
  \end{itemize}
\end{practiceans}

% ---- Q31
\Needspace{9\baselineskip}
\begin{practiceans}
  \correct{A}
  The field is into the page and the conductor lies in the page, so the conductor is at \ang{90} to the
  field, whatever its angle within the page. $F = BIl = (0.25)(3.0)l$.
  \begin{itemize}
    \item \xmark\ B, C and D: the \ang{40} is between the wire and the horizontal, not between the wire
          and the field.
  \end{itemize}
\end{practiceans}

% ---- Q32
\Needspace{9\baselineskip}
\begin{practiceans}
  \correct{C}
  $Y$ hangs below $X$, so the magnetic force on $Y$ is upwards and equals the weight of $Y$: the rods
  attract, so the currents are in the same direction. By Newton's third law, $Y$ pulls $X$ down with an
  equal force. The support must hold the weight of $X$ plus this force, which equals the weight of $Y$.
  \begin{itemize}
    \item \xmark\ A: opposite currents would repel. \quad \xmark\ B: the force on $X$ is towards $Y$,
          which is downwards.
    \item \xmark\ D: by Newton's third law the two forces are equal in size.
  \end{itemize}
\end{practiceans}

% ---- Q33
\Needspace{9\baselineskip}
\begin{practiceans}
  \correct{B}
  The magnet pushes the rod up with a force $BIl$. By Newton's third law, the rod pushes the magnet down
  with $BIl$. The balance shows this extra force as extra mass $\frac{BIl}{g}$, so the reading is
  $0.300 + \frac{BIl}{9.8}$.
  \begin{itemize}
    \item \xmark\ A: adds a force (newtons) to a mass (kilograms).
    \item \xmark\ C and D: the magnet is pushed down, not up.
  \end{itemize}
\end{practiceans}

% ---- Q34
\Needspace{9\baselineskip}
\begin{practiceans}
  \correct{D}
  The currents are always in the same direction, so the force is always attraction. Its size depends on
  $I_1 I_2 = I^2$, which rises and falls to zero twice in each cycle but never becomes negative: a series
  of pulses of attraction.
  \begin{itemize}
    \item \xmark\ A: the current changes all the time, so the force does too.
    \item \xmark\ B: currents in the same direction attract.
    \item \xmark\ C: both currents reverse together, so the force never becomes repulsion.
  \end{itemize}
\end{practiceans}

% ---- Q35
\Needspace{9\baselineskip}
\begin{practiceans}
  \correct{A}
  By Newton's third law, the forces the wires exert on each other are equal in size and opposite in
  direction. ($\frac{F}{l} = \frac{\mu_0 I_1 I_2}{2\pi r}$ uses the product $I_1 I_2$, the same for both
  wires.)
  \begin{itemize}
    \item \xmark\ B and C: the forces are equal, whichever wire has the larger current.
    \item \xmark\ D: the forces are in opposite directions.
  \end{itemize}
\end{practiceans}

% ---- Q36
\Needspace{9\baselineskip}
\begin{practiceans}
  \correct{D}
  The currents are opposite, so the wires repel: the force on $X$ is upwards (away from $Y$). The force is
  proportional to $I_X I_Y$, so tripling $I_X$ makes it $3F$. The answer is a long arrow (three times $F$)
  pointing up.
  \begin{itemize}
    \item \xmark\ A and B: these arrows are only the size of $F$.
    \item \xmark\ C: the force on $X$ points away from $Y$, which is upwards.
  \end{itemize}
\end{practiceans}

% =====================================================================
\anstopic{Topic 3 - Electromagnetic Induction}

% ---- Q37
\begin{practiceans}
  \correct{B}
  The field is at \ang{30} to the plane, so it is at \ang{60} to the normal (the line at right angles to
  the loop). $\Phi = BA\cos 60^\circ = 0.40 \times 0.050 \times 0.5 = \qty{0.010}{Wb}$.
  \begin{itemize}
    \item \xmark\ A: some of the field passes through the loop.
    \item \xmark\ C: uses \ang{30} as the angle to the normal. \quad \xmark\ D: ignores the angle.
  \end{itemize}
\end{practiceans}

% ---- Q38
\Needspace{9\baselineskip}
\begin{practiceans}
  \correct{D}
  $\Phi = BA\cos\theta$ depends on the field strength, the area and the angle. Moving the loop sideways
  inside a uniform field changes none of these, so the flux stays the same (and no emf is induced).
  \begin{itemize}
    \item \xmark\ A, B and C: each of these changes $B$, $\theta$ or $A$.
  \end{itemize}
\end{practiceans}

% ---- Q39
\Needspace{9\baselineskip}
\begin{practiceans}
  \correct{C}
  $\Delta\Phi = BA = 0.60 \times \num{4.0e-3} = \qty{2.4e-3}{Wb}$.
  $\varepsilon = N\dfrac{\Delta\Phi}{\Delta t} = 250 \times \dfrac{\num{2.4e-3}}{0.050} = \qty{12}{V}$.
  \begin{itemize}
    \item \xmark\ A: uses one turn. \quad \xmark\ B: forgets to divide by $\Delta t$.
    \item \xmark\ D: forgets the area.
  \end{itemize}
\end{practiceans}

% ---- Q40
\Needspace{9\baselineskip}
\begin{practiceans}
  \correct{A}
  $\varepsilon = -N\frac{\Delta\Phi}{\Delta t}$ is proportional to minus the gradient of the flux graph.
  0--2 s: positive gradient, so a constant negative emf. 2--3 s: zero gradient, so zero emf. 3--4 s:
  negative gradient twice as steep, so a positive emf twice as large.
  \begin{itemize}
    \item \xmark\ B: the signs are reversed. \quad \xmark\ C: this copies the flux graph.
    \item \xmark\ D: the two emfs should be different sizes.
  \end{itemize}
\end{practiceans}

% ---- Q41
\Needspace{9\baselineskip}
\begin{practiceans}
  \correct{D}
  The field of a long wire gets weaker with distance ($B \propto \frac{1}{r}$), and it changes less and
  less quickly further away. Moving at constant speed, the loop's flux changes quickly at first and then
  more slowly, so the emf decreases along a curve towards zero.
  \begin{itemize}
    \item \xmark\ A: the rate of change of flux is not constant.
    \item \xmark\ B: the decrease is not linear, and the emf does not reach zero at a definite time.
    \item \xmark\ C: the emf decreases, not increases.
  \end{itemize}
\end{practiceans}

% ---- Q42
\Needspace{9\baselineskip}
\begin{practiceans}
  \correct{D}
  Lenz's law: the induced current opposes the change that causes it, here the magnet moving away. So the
  coil attracts the magnet. To attract a north pole, the right-hand end of the coil acts as a south pole.
  \begin{itemize}
    \item \xmark\ A and C: repulsion would help the magnet move away, creating energy from nothing.
    \item \xmark\ B: a north pole would repel the magnet's north pole.
  \end{itemize}
\end{practiceans}

% ---- Q43
\Needspace{9\baselineskip}
\begin{practiceans}
  \correct{C}
  An emf is induced in the second coil only while the flux changes. With a steady current, the flux is
  constant, so the reading is zero. When the switch opens, the flux decreases (instead of increasing), so
  the induced current is in the opposite direction: a brief flick to the right.
  \begin{itemize}
    \item \xmark\ A and D: a steady current gives no induced current.
    \item \xmark\ B: a decreasing flux gives the opposite direction to an increasing flux.
  \end{itemize}
\end{practiceans}

% ---- Q44
\Needspace{9\baselineskip}
\begin{practiceans}
  \correct{D}
  While entering, the flux through the loop increases, so an emf is induced and the lamp lights. While
  completely inside the uniform field, the flux is constant: no emf. While leaving, the flux decreases:
  an emf in the opposite direction, so the lamp lights again.
  \begin{itemize}
    \item \xmark\ A: there is no emf while the flux is constant. \quad \xmark\ B: this is when the flux is
          constant.
    \item \xmark\ C: constant speed does not matter; a changing flux does.
  \end{itemize}
\end{practiceans}

% ---- Q45
\Needspace{9\baselineskip}
\begin{practiceans}
  \correct{A}
  Arrangement 1: the rod cuts across the field lines at a steady rate, so it sweeps flux at a constant
  rate: a constant emf. Arrangement 2: the rod moves within the plane that contains the field lines, so it
  never cuts across them: no emf.
  \begin{itemize}
    \item \xmark\ B: swapped. \quad \xmark\ C: no field lines are cut in arrangement 2.
    \item \xmark\ D: the rate of rotation is constant, so the emf is constant.
  \end{itemize}
\end{practiceans}

% ---- Q46
\Needspace{9\baselineskip}
\begin{practiceans}
  \correct{B}
  The area of the circuit grows by $lv$ each second, so $\varepsilon = \dfrac{\Delta\Phi}{\Delta t} = Blv
  = 0.30 \times 0.50 \times 2.0 = \qty{0.30}{V}$. $I = \dfrac{\varepsilon}{R} = \dfrac{0.30}{1.5} = \qty{0.20}{A}$.
  \begin{itemize}
    \item \xmark\ A: forgets $v$. \quad \xmark\ C: forgets $l$.
    \item \xmark\ D: multiplies by $R$ instead of dividing.
  \end{itemize}
\end{practiceans}

% ---- Q47
\Needspace{9\baselineskip}
\begin{practiceans}
  \correct{C}
  The falling magnet changes the flux through the copper tube, inducing eddy currents. By Lenz's law
  they oppose the magnet's motion, so it falls more slowly. Plastic is an insulator, so no currents flow
  in it.
  \begin{itemize}
    \item \xmark\ A: copper is not magnetic, but it is a good conductor.
    \item \xmark\ B: speeding the magnet up would create energy from nothing.
    \item \xmark\ D: copper is not attracted to a magnet.
  \end{itemize}
\end{practiceans}

% ---- Q48
\Needspace{9\baselineskip}
\begin{practiceans}
  \correct{B}
  $\dfrac{V_\mathrm{p}}{V_\mathrm{s}} = \dfrac{N_\mathrm{p}}{N_\mathrm{s}}$, so
  $N_\mathrm{s} = 1200 \times \dfrac{12}{240} = 60$.
  \begin{itemize}
    \item \xmark\ A: this is not a number of turns. \quad \xmark\ C: doubles $N_\mathrm{p}$.
    \item \xmark\ D: the ratio is upside down.
  \end{itemize}
\end{practiceans}

% ---- Q49
\Needspace{9\baselineskip}
\begin{practiceans}
  \correct{C}
  $V_\mathrm{p}I_\mathrm{p} = V_\mathrm{s}I_\mathrm{s}$. If $I_\mathrm{s} = 5I_\mathrm{p}$, then
  $V_\mathrm{s} = \frac{V_\mathrm{p}}{5}$: a step-down transformer, with
  $N_\mathrm{p} : N_\mathrm{s} = V_\mathrm{p} : V_\mathrm{s} = 5 : 1$.
  \begin{itemize}
    \item \xmark\ A and B: a step-up transformer makes the secondary current smaller.
    \item \xmark\ D: the ratio is upside down.
  \end{itemize}
\end{practiceans}

% ---- Q50
\Needspace{9\baselineskip}
\begin{practiceans}
  \correct{B}
  $V_\mathrm{s} = V_\mathrm{p}\dfrac{N_\mathrm{s}}{N_\mathrm{p}}$. Halving $N_\mathrm{p}$ doubles $V_\mathrm{s}$.
  \begin{itemize}
    \item \xmark\ A and D: these halve $V_\mathrm{s}$. \quad \xmark\ C: the ratio, and $V_\mathrm{s}$, stay the same.
  \end{itemize}
\end{practiceans}

% ---- Q51
\Needspace{9\baselineskip}
\begin{practiceans}
  \correct{D}
  Stepping up to a high voltage makes the current small ($P = VI$), so the heating loss $I^2R$ in the
  lines is small. Near homes the voltage is stepped down to a safe level.
  \begin{itemize}
    \item \xmark\ A and C: a large current would cause large $I^2R$ losses.
    \item \xmark\ B: transmitting at a low voltage would waste far more energy.
  \end{itemize}
\end{practiceans}

% ---- Q52
\Needspace{9\baselineskip}
\begin{practiceans}
  \correct{B}
  $I = \dfrac{P}{V} = \dfrac{\num{50e6}}{\num{250e3}} = \qty{200}{A}$. Power loss
  $= I^2R = 200^2 \times 5.0 = \qty{2.0e5}{W}$.
  \begin{itemize}
    \item \xmark\ A: this is $IR$, the voltage drop in the line (in volts), not a power.
    \item \xmark\ C: this would be the loss at \qty{25}{kV}.
    \item \xmark\ D: $\frac{V^2}{R}$ needs the voltage \emph{across the line}, which is not \qty{250}{kV}.
  \end{itemize}
\end{practiceans}

% ---- Q53
\Needspace{9\baselineskip}
\begin{practiceans}
  \correct{A}
  The changing flux induces eddy currents in the iron core, which heat it. Thin insulated sheets break
  up the current paths, so the eddy currents are much smaller and less energy is wasted.
  \begin{itemize}
    \item \xmark\ B: a larger coil resistance would waste more energy.
    \item \xmark\ C: laminations do not stop flux leakage. \quad \xmark\ D: this is not the reason.
  \end{itemize}
\end{practiceans}

% ---- Q54
\Needspace{9\baselineskip}
\begin{practiceans}
  \correct{A}
  A steady DC current gives a constant flux. With no change in flux, no emf is induced in the secondary
  coil, so the output is \qty{0}{V}. (Transformers work only with a changing current, such as AC.)
  \begin{itemize}
    \item \xmark\ B, C and D: these assume a changing flux.
  \end{itemize}
\end{practiceans}

% =====================================================================
\anstopic{Topic 4 - Applications of the Motor Effect}

% ---- Q55
\begin{practiceans}
  \correct{D}
  $\tau = nIAB\sin\theta$, where $\theta$ is the angle between the field and the normal to the coil. The
  plane is parallel to the field, so $\theta = 90^\circ$ and the torque is a maximum:
  $\tau = 100 \times 2.0 \times 0.010 \times 0.20 = \qty{0.40}{N.m}$.
  \begin{itemize}
    \item \xmark\ A: the torque is zero when the plane is at right angles to the field, not parallel.
    \item \xmark\ B: forgets the 100 turns. \quad \xmark\ C: half the correct value.
  \end{itemize}
\end{practiceans}

% ---- Q56
\Needspace{9\baselineskip}
\begin{practiceans}
  \correct{B}
  Each long side stays at right angles to the field, so the force $F = lIB$ is unchanged. But as the coil
  turns, the forces act at a smaller perpendicular distance from the axis, so the torque decreases. (With
  the plane at \ang{60} to the field, the normal is at \ang{30}, so $\tau = \tau_\mathrm{max}\sin 30^\circ$,
  half the maximum.)
  \begin{itemize}
    \item \xmark\ A and D: the force on each side does not change.
    \item \xmark\ C: the torque decreases.
  \end{itemize}
\end{practiceans}

% ---- Q57
\Needspace{9\baselineskip}
\begin{practiceans}
  \correct{C}
  Every half turn, the split ring swaps the connections to the supply, reversing the current in the
  coil. This keeps the torque in the same direction, so the motor keeps turning the same way.
  \begin{itemize}
    \item \xmark\ A: a DC motor runs on DC. \quad \xmark\ B: the magnets provide the field.
    \item \xmark\ D: this is not its purpose.
  \end{itemize}
\end{practiceans}

% ---- Q58
\Needspace{9\baselineskip}
\begin{practiceans}
  \correct{C}
  The back emf is induced by the coil turning in the field. It is proportional to the rate of change of
  flux, and so to the angular velocity: a straight line through the origin.
  \begin{itemize}
    \item \xmark\ A: the back emf depends on speed. \quad \xmark\ B: a faster motor has a \emph{larger}
          back emf.
    \item \xmark\ D: the relationship is linear, not curved.
  \end{itemize}
\end{practiceans}

% ---- Q59
\Needspace{9\baselineskip}
\begin{practiceans}
  \correct{A}
  A slower motor has a smaller back emf. The net voltage across the coil ($V - \varepsilon$) increases,
  so the current increases. (This larger current gives the larger torque needed for the load.)
  \begin{itemize}
    \item \xmark\ B: a smaller back emf means a larger current.
    \item \xmark\ C and D: the back emf falls when the motor slows.
  \end{itemize}
\end{practiceans}

% ---- Q60
\Needspace{9\baselineskip}
\begin{practiceans}
  \correct{B}
  $V = \varepsilon_\mathrm{back} + IR$, so $\varepsilon_\mathrm{back} = 24 - 1.5 \times 2.0 = \qty{21}{V}$.
  \begin{itemize}
    \item \xmark\ A: this is $IR$. \quad \xmark\ C: this is the supply voltage.
    \item \xmark\ D: adds $IR$ instead of subtracting.
  \end{itemize}
\end{practiceans}

% ---- Q61
\Needspace{9\baselineskip}
\begin{practiceans}
  \correct{A}
  In an AC induction motor the rotor current is induced by the rotating field of the stator. There is
  no electrical connection to the rotor, so no brushes or commutator are needed.
  \begin{itemize}
    \item \xmark\ B: the stator coils make the rotating field. \quad \xmark\ C: the rotor is the part that
          turns.
    \item \xmark\ D: the AC supply drives the stator coils.
  \end{itemize}
\end{practiceans}

% ---- Q62
\Needspace{9\baselineskip}
\begin{practiceans}
  \correct{C}
  A current is induced in the rotor only if the field moves relative to it. At the same speed as the
  field, the flux through the rotor would not change, so there would be no current and no torque. The
  rotor always lags a little behind (slip).
  \begin{itemize}
    \item \xmark\ A, B and D: these are not the reason the rotor must turn more slowly.
  \end{itemize}
\end{practiceans}

% ---- Q63
\Needspace{9\baselineskip}
\begin{practiceans}
  \correct{D}
  The emf in the coil alternates, but the split ring swaps the connections every half turn, so the output
  never changes sign: a fluctuating DC, shaped like the AC curve with the negative halves flipped up.
  \begin{itemize}
    \item \xmark\ A: this is the output with slip rings (an AC generator).
    \item \xmark\ B: the emf changes smoothly, not suddenly. \quad \xmark\ C: the output is not steady.
  \end{itemize}
\end{practiceans}

% ---- Q64
\Needspace{9\baselineskip}
\begin{practiceans}
  \correct{A}
  Turning twice as fast makes the flux change twice as quickly, so the peak emf doubles. There are also
  twice as many turns each second, so the frequency doubles.
  \begin{itemize}
    \item \xmark\ B, C and D: both quantities double.
  \end{itemize}
\end{practiceans}

% ---- Q65
\Needspace{9\baselineskip}
\begin{practiceans}
  \correct{D}
  Two bulbs in parallel have the lowest resistance, so the largest current flows. The current in the
  coil produces a force that opposes the turning (Lenz's law); a larger current means a larger opposing
  force. In energy terms, more electrical power out needs more work in.
  \begin{itemize}
    \item \xmark\ A: with no current there is no opposing force, so it is easiest to turn.
    \item \xmark\ B: bulbs in series have a high resistance, so the current is small.
    \item \xmark\ C: one bulb draws less current than two in parallel.
  \end{itemize}
\end{practiceans}

% ---- Q66
\Needspace{9\baselineskip}
\begin{practiceans}
  \correct{C}
  The moving magnets induce eddy currents in the rails, which oppose the motion (Lenz's law) and turn
  kinetic energy into heat. The induced currents depend on the relative speed, so the braking force
  decreases as the train slows.
  \begin{itemize}
    \item \xmark\ A: the force is greatest at high speed. \quad \xmark\ B: no contact is needed.
    \item \xmark\ D: the force falls to zero as the speed falls to zero, so friction brakes are needed to
          stop completely.
  \end{itemize}
\end{practiceans}

% ---- Q67
\Needspace{9\baselineskip}
\begin{practiceans}
  \correct{B}
  If an induced current helped the change that caused it, it would make the change bigger, which would
  make a bigger current, and so on, giving more and more energy for nothing. Energy is conserved, so the
  induced current must oppose the change: work must be done to make the change.
  \begin{itemize}
    \item \xmark\ A, C and D: these are not what Lenz's law says.
  \end{itemize}
\end{practiceans}

% ---- Q68
\Needspace{9\baselineskip}
\begin{practiceans}
  \correct{A}
  The slots break up the paths of the eddy currents, so the currents and their braking force are much
  smaller. The slotted plate loses its energy more slowly and swings for longer.
  \begin{itemize}
    \item \xmark\ B: the braking depends on the eddy currents, not the mass.
    \item \xmark\ C: the eddy currents are different. \quad \xmark\ D: the slots reduce the eddy currents.
  \end{itemize}
\end{practiceans}

% ---- Q69
\Needspace{9\baselineskip}
\begin{practiceans}
  \correct{B}
  Input power $= VI = 12 \times 2.5 = \qty{30}{W}$. Useful output power $= mgv = 1.5 \times 9.8 \times
  1.2 = \qty{17.6}{W}$. Efficiency $= \dfrac{17.6}{30} = 59\%$.
  \begin{itemize}
    \item \xmark\ A: this is the fraction wasted. \quad \xmark\ C: ignores the energy lost as heat.
    \item \xmark\ D: the fraction is upside down; efficiency cannot be over 100\%.
  \end{itemize}
\end{practiceans}

% ---- Q70
\Needspace{9\baselineskip}
\begin{practiceans}
  \correct{C}
  $\Delta\Phi = BA - 0 = 0.40 \times 0.020 = \qty{8.0e-3}{Wb}$.
  $\varepsilon_\mathrm{av} = N\dfrac{\Delta\Phi}{\Delta t} = 50 \times \dfrac{\num{8.0e-3}}{0.025} = \qty{16}{V}$.
  \begin{itemize}
    \item \xmark\ A: uses one turn. \quad \xmark\ B: uses twice the time.
    \item \xmark\ D: forgets the area.
  \end{itemize}
\end{practiceans}

% ---- Q71
\Needspace{9\baselineskip}
\begin{practiceans}
  \correct{D}
  A generator turns mechanical energy into electrical energy. A motor's coil also turns in a magnetic
  field, so it acts as a generator too; the emf it produces opposes the supply (back emf).
  \begin{itemize}
    \item \xmark\ A: a motor turns electrical energy into mechanical energy.
    \item \xmark\ B: a generator only needs a coil moving in a magnetic field.
    \item \xmark\ C: they use the same parts (coil, magnets, commutator or slip rings).
  \end{itemize}
\end{practiceans}

% ---- Q72
\Needspace{9\baselineskip}
\begin{practiceans}
  \correct{B}
  The moving magnet changes the flux in the copper, inducing eddy currents. By Lenz's law, the force
  between the currents and the magnet opposes their \emph{relative} motion, so the block is dragged
  along in the same direction as the magnet.
  \begin{itemize}
    \item \xmark\ A: copper is not magnetic, but it conducts, so eddy currents flow.
    \item \xmark\ C: a force opposite to the magnet's motion would increase the relative motion.
    \item \xmark\ D: copper is not attracted to a magnet.
  \end{itemize}
\end{practiceans}

\end{practiceanswers}
\end{document}
